\chapter{Literature Review and Research Gaps}
\label{Chapter2}
\lhead{Chapter 2. \emph{Literature Review and Gaps}}

Prior work on container runtime security falls into four groups: production
rule engines, anomaly-based syscall detectors, in-kernel enforcers, and
cross-container research. Each group solves a real problem, and each leaves
a specific gap that CausalTrace is engineered to close. We frame the review
around those gaps.

%─────────────────────────────────────────────────────────────────────────────
\section{Production runtime security tools}
\label{sec:lit-production}
%─────────────────────────────────────────────────────────────────────────────

Three open-source tools dominate the production market. Falco~\cite{falco2024}
matches kernel events against a YAML rule library through a modern-eBPF
probe. Tetragon~\cite{tetragon2024}, a sub-project of
Cilium~\cite{cilium2024}, compiles \emph{TracingPolicies} to eBPF kprobes and
supports in-kernel enforcement through signal delivery and socket override.
Tracee~\cite{tracee2024} focuses on forensics: it emits events for post-hoc
analysis rather than enforcing in the hot path.

All three operate on per-container streams. Rules match inside one
container's own syscall trace or against a deny-list; there is no primitive
that expresses \emph{container $A$ spawns a shell within five seconds of
container $B$ connecting to a previously-unseen port}. A correlation layer
can be built on top of any of them, but the round-trip to userspace that
the layer forces up wipes out the latency advantage that in-kernel detection
was supposed to provide.

Ghimire~\emph{et al.}~\cite{ghimire2025patrol} (eBPF-PATROL) are the closest
published baseline to our Tier~1 design. They intercept six dangerous
syscalls (\texttt{clone}, \texttt{execve}, \texttt{openat}, \texttt{setns},
\texttt{unshare}, \texttt{ptrace}) and enforce through BPF LSM hooks at
$\sim$100\,ns and \texttt{bpf\_override\_return} at $\sim$200\,ns. They
miss the same cross-container class that Falco and Tetragon miss, and they
require a twelve-entry infrastructure whitelist (documented by them and
reproduced in our own mid-review) to keep Docker's own runtime from
tripping the rules.

%─────────────────────────────────────────────────────────────────────────────
\section{Anomaly-based syscall detectors}
\label{sec:lit-bosc}
%─────────────────────────────────────────────────────────────────────────────

The Bag-of-System-Calls (BoSC) family represents a container's behaviour as
a frequency vector over syscall opcodes. Bertinatto~\emph{et
al.}~\cite{bertinatto2024} attach a raw tracepoint to \texttt{sys\_enter},
filter by mount namespace, and build per-process BoSC vectors. The method
cleanly catches one hardcoded container-escape pattern
(\texttt{openat("/proc/1/ns/mnt")} followed by
\texttt{setns(CLONE\_NEWNS)}) and misses every other class of attack.

Autoencoder-based detectors~\cite{kotenko2024,desilva2022micro,xia2024selfsup}
and sequence models~\cite{chen2021lstm,fournier2023lm} lift this work by
learning a density estimate over benign traces. They share the structural
weakness: the model learns the \emph{marginal} distribution per container
but not the joint structure between containers, so an attacker whose
individual syscall frequencies stay close to the marginal can compose an
attack whose joint signature is far from any single container's
calibrated regime.

%─────────────────────────────────────────────────────────────────────────────
\section{In-kernel enforcement and adversarial evasion}
\label{sec:lit-enforcement}
%─────────────────────────────────────────────────────────────────────────────

A smaller body of work targets enforcement in kernel space. Zhang~\emph{et
al.}~\cite{zhang2024rtids} build a neural-network inference engine that
runs inside eBPF through tail calls and integer arithmetic, reporting F1
scores of 0.933 and 0.992 on packet-level intrusion datasets with 5\,KB of
in-kernel memory. The design is elegant at the verifier boundary, but the
detection target is packet flows, not multi-container chains.

Takabayashi and Mambo~\cite{takabayashi2024evasion} study the adversarial
other side: how an attacker can \emph{evade} eBPF-based detectors by
\texttt{LD\_PRELOAD}-rewriting syscall arguments, hiding process files
through \texttt{getdents} rewriting, or dropping payload packets through
XDP. Their evasion results motivate the strict-invariant design
philosophy of CausalTrace's Tier~1: we hook on physical contract violations
(a process whose standard input points to a socket is a reverse shell) that
the attacker cannot semantically rewrite away.

%─────────────────────────────────────────────────────────────────────────────
\section{Cross-container attack research}
\label{sec:lit-cross}
%─────────────────────────────────────────────────────────────────────────────

He~\emph{et al.}~\cite{he2023cross} document the cross-container attack
class. They catalogue five vulnerable online services in which offensive
eBPF features break container isolation, and three cloud providers' default
Kubernetes clusters vulnerable to cross-node attacks. Their contribution is
attack discovery; they propose no defender.

Zhang~\emph{et al.}~\cite{zhang2025path} (Pamir-Detector) monitors
path-resolution syscalls for filesystem-isolation bypasses (symlink attacks,
procfs traversal, host-process-mediated file access). Their coverage is
one escape family; CausalTrace's Tier~1 handlers cover complementary
families (socket-stdio reverse shells, namespace manipulation, setuid-0
from non-root). No published defender models the joint behaviour of
multiple containers mathematically and detects composite attacks that never
trip a single-container rule. That is the gap this thesis closes.

%─────────────────────────────────────────────────────────────────────────────
\section{Mathematical foundations}
\label{sec:lit-math}
%─────────────────────────────────────────────────────────────────────────────

The Tier~3 detector reuses four classical constructions as stated; no new
mathematical theorems about them are claimed.

Canonical-correlation analysis~\cite{hotelling1936} supplies the restriction
maps $F_u, F_v$ per edge of the communication graph. The Mahalanobis
distance~\cite{mahalanobis1936} scores residuals against the learned
coupling covariance. The Count-Min Sketch~\cite{cormode2005cms} gives our
in-kernel bigram counter at $4 \times 128$ per container. Cellular-sheaf
spectral theory~\cite{hansen2019spectral,curry2014sheaves,robinson2014topological}
provides the sheaf Laplacian~$L_F$ and the Rayleigh-quotient interpretation
as a coupling-consistency measure. The sheaf specialises to the ordinary
graph Laplacian of spectral graph theory~\cite{chung1997spectral} whenever
every restriction map is the identity.

%─────────────────────────────────────────────────────────────────────────────
\section{Capability matrix and research gaps}
\label{sec:gaps}
%─────────────────────────────────────────────────────────────────────────────

Table~\ref{tab:capability-matrix} summarises where each reviewed system sits
against the capabilities we consider essential for defending a composite
cross-container attack. A single row is filled in by CausalTrace.

\begin{table}[t]
\centering
\caption[Capability matrix of reviewed runtime security systems]
{Capabilities of reviewed runtime security systems. \ding{51}\,=\,full,
$\sim$\,=\,partial, \ding{55}\,=\,absent. $^{\dagger}$In-kernel enforcement
is possible via \texttt{bpf\_send\_signal} but is community-discouraged in
production because it terminates the entire container process.}
\label{tab:capability-matrix}
\footnotesize
\setlength\tabcolsep{3.5pt}
\begin{tabular}{@{}l c c c c c c@{}}
\toprule
System & Cross- & Kernel- & Sub-$\mu$s & OOD & No labelled & Attacker-IP \\
       & container & native enf. & latency & detect. & train data & severance \\
\midrule
Falco~\cite{falco2024}          & \ding{55} & \ding{55}$^{\dagger}$ & \ding{55} & \ding{55} & \ding{51} & \ding{55} \\
Tetragon~\cite{tetragon2024}    & \ding{55} & $\sim$    & \ding{51} & \ding{55} & \ding{51} & \ding{55} \\
Tracee~\cite{tracee2024}        & \ding{55} & \ding{55} & \ding{55} & \ding{55} & \ding{51} & \ding{55} \\
BoSC~\cite{bertinatto2024}      & \ding{55} & \ding{55} & \ding{55} & \ding{55} & \ding{55} & \ding{55} \\
eBPF-PATROL~\cite{ghimire2025patrol}
                                & \ding{55} & \ding{51} & \ding{51} & \ding{55} & \ding{51} & \ding{55} \\
LSTM~/~LM~\cite{chen2021lstm,fournier2023lm}
                                & \ding{55} & \ding{55} & \ding{55} & $\sim$    & \ding{55} & \ding{55} \\
Pamir-Detector~\cite{zhang2025path}
                                & \ding{55} & \ding{51} & \ding{51} & \ding{55} & \ding{51} & \ding{55} \\
\textbf{CausalTrace (this work)}
                                & \textbf{\ding{51}} & \textbf{\ding{51}}
                                & \textbf{\ding{51}} & \textbf{\ding{51}}
                                & \textbf{\ding{51}} & \textbf{\ding{51}} \\
\bottomrule
\end{tabular}
\end{table}

From this matrix and the preceding review we distil five research gaps.

\paragraph{G1 --- Semantic blindness of rule-based detectors.}
Baseline~B (our mid-review Modified PATROL) detected all five attack classes
but required twelve infrastructure whitelist entries
(\texttt{apt-get}, \texttt{gpgv}, \texttt{runc}, \texttt{containerd}, and
eight others) to keep normal Docker operations from tripping the rules. Each
entry is both a maintenance burden and an evasion vector. The gap is
semantic: identical syscalls have different intent depending on context,
and no amount of rule tuning closes a gap that lives in meaning rather than
syntax.

\paragraph{G2 --- No cross-container chain detection.}
Every reviewed defender operates on per-container events. Multi-stage
chains---Log4Shell in a webapp, Dirty-Pipe for privilege escalation, a
cgroup escape to the host---are invisible to a single-event monitor because
each stage is locally plausible. He~\emph{et al.}~\cite{he2023cross}
document the attack class and propose no defender.

\paragraph{G3 --- Userspace enforcement bottleneck.}
PATROL-class architectures route every decision through a userspace Go
controller: eBPF probe $\rightarrow$ ring buffer $\rightarrow$ controller
$\rightarrow$ \texttt{kill()}. Mid-review measurements on that path
recorded approximately $23\,\mu$s round-trip; the triggering syscall
completes before the kill arrives.

\paragraph{G4 --- Coarse enforcement target.}
Kernel-native enforcement delivered through \texttt{bpf\_send\_signal(SIGKILL)}
targets the container process itself. In a deployment serving legitimate
users alongside the attacker, this is collateral damage---every in-flight
request terminates the moment an invariant fires. No reviewed system
severs at attacker-IP granularity.

\paragraph{G5 --- No mathematical model of inter-container coupling.}
The object that encodes \emph{which containers may communicate, at what
cadence, and with what joint syscall distribution} is absent from every
reviewed system. BoSC operates on per-container marginals; network
detectors operate on per-flow features; neither captures the joint
structure that a sheaf Laplacian makes explicit.

Chapters~\ref{Chapter3}--\ref{Chapter6} close these five gaps, respectively
through the math, the architecture, the implementation, and the evaluation.
